% !TeX root = ../main.tex
\section{An Approach for Guitar Virtual Amp}
\label{sec:application}

This section derives a design from the review for the application described
in Section~\ref{sec:intro}. It is a proposal: none of it has been measured
yet, and Section~\ref{sec:protocol} gives the protocol to measure it.

\subsection{What the application already has}

\begin{itemize}
\item \emph{A guitar stem per part.} Separation with a BS-RoFormer six-stem
model~\cite{lu2023bsroformer} (Demucs as fallback~\cite{rouard2023htdemucs}),
and a split of the stereo guitar stem by pan position into parts that play
different things~\cite{deangelo2026tabreview}.
\item \emph{The notes and the grid.} Onsets, offsets, pitches and techniques
of every guitar part, beats, bars and tempo from Beat This!~\cite{foscarin2024beatthis};
for songs whose Guitar Pro tab the user owns, a tab aligned to the recording.
\item \emph{An offline renderer of the real chain.} The same pedal code as the
live path, at 48\,kHz in 256-sample blocks, already used to render training
data through presets. The full chain of 17 pedals costs about 430\,ns per
sample on one core, and a NAM WaveNet about 1.9\,\textmu s.
\item \emph{Presets with song metadata.} 20 factory presets, user presets with
\texttt{song} and \texttt{artist} fields, and a tone suggestion on the Tabs
page that prefers a user preset tagged with the song; the current fallback
guesses a factory preset from the track's instrument and name.
\item \emph{Capture and cabinet players.} A NAM pedal and an impulse-response
convolver, with TONE3000 browsing under its API terms.
\item \emph{Data and compute.} EGDB, Guitar-TECHS, GuitarSet, GOAT, IDMT,
EGSet12, MUSDB18 and MoisesDB on disk, a mixing script that places labelled
guitar over backing, and DGX Spark nodes for training.
\end{itemize}

\subsection{Constraints}

The chain is mono, works on 16-bit samples, and is written in Rust without
gradients, so the estimator must treat it as a black box. Stereo effects
cannot be reproduced. The estimate may run offline on a desktop (Python
sidecars are acceptable, as for separation), and its output, a preset, also
works on Android. The analysed audio must be the user's own: Spotify's
developer policy forbids analysing its content, so the feature applies to the
same local files and downloads as transcription. Models that ship must be
trained on data that allow it (Section~\ref{sec:datasets}), and capture
libraries can be searched by metadata but not bulk downloaded. Patent
CN118865922B claims a pipeline of the same shape~\cite{huanyin2025presetpatent};
its claims should be reviewed before distributing such a feature where it
applies.

\subsection{Proposed pipeline}

Fig.~\ref{fig:pipeline} shows the pipeline. It follows the consensus of
Sections~\ref{sec:recognition}--\ref{sec:products}: predict, refine by
rendering the real chain, regularise with a preset prior, and finish with a
linear match.

\begin{figure}[t]
\centering
\begin{tikzpicture}[
  node distance=1.8mm,
  box/.style={draw, rounded corners=1pt, font=\scriptsize, align=center,
              minimum height=5mm, inner sep=2pt, fill=white, text width=36mm},
  have/.style={box, fill=gray!15},
  arr/.style={-{Stealth[length=1.6mm]}, semithick}]
  \node[have] (sep) {1. Separation, pan split, notes, beats};
  \node[box, below=of sep] (seg) {2. Excerpts per part and section};
  \node[box, below=of seg] (di) {3. Pseudo-dry signal: resynthesis\\from notes, remover, or player DI};
  \node[box, below=of di] (init) {4. Initial chains: classifier,\\regressor, preset retrieval (top-$k$)};
  \node[box, below=of init] (ref) {5. CMA-ES through the offline\\renderer; spectral, embedding, prior};
  \node[box, below=of ref] (time) {6. Delay, reverb, modulation\\from notes and tempo};
  \node[box, below=of time] (lin) {7. Linear match: EQ pedal or IR};
  \node[have, below=of lin] (out) {8. Preset tagged with the song};
  \draw[arr] (sep) -- (seg);
  \draw[arr] (seg) -- (di);
  \draw[arr] (di) -- (init);
  \draw[arr] (init) -- (ref);
  \draw[arr] (ref) -- (time);
  \draw[arr] (time) -- (lin);
  \draw[arr] (lin) -- (out);
\end{tikzpicture}
\caption{Proposed tone estimation pipeline. Grey boxes exist in the
application today.}
\label{fig:pipeline}
\end{figure}

\paragraph{Excerpts} For each guitar part, pick excerpts of 5--20\,s where the
part is loud relative to the residual of the mix, other guitars are quiet, and
the notes cover several registers, using the transcription to measure all
three. Cluster excerpts by an effect embedding to find tone changes between
sections, and estimate one preset per cluster.

\paragraph{Pseudo-dry signals} Three sources, used together when available.
(a) \emph{Resynthesis from the notes}: render the transcribed (or aligned tab)
notes as a clean electric guitar. The result is aligned with the stem by
construction, but its timbre is not the original guitar's; the first version
can use the existing SoundFont synthesizer, and a learned DI renderer in the
manner of GuitarFlow~\cite{loth2025guitarflow}, trained on DI corpora that
allow it, would narrow the gap. This route is new as far as we found. (b)
\emph{Removal}: a model trained to map renders through the application's own
chains, and through EGDB-PG presets, back to their DI~\cite{lee2024distortion,yeh2026egvae}.
(c) \emph{The player's DI}: the user plays along in practice mode, which
requires tapping the signal before the chain. (c) has the right guitar and
wrong timing; (a) the right timing and wrong guitar; (b) both, with removal
errors.

\paragraph{Template chain} Search over one fixed chain: noise gate,
compressor, drive, amplifier or NAM capture, cabinet, equaliser, one
modulation effect, delay, reverb, each with an on/off switch, about 35
continuous parameters plus the amplifier voicing, the impulse response and the
capture. A fixed order follows the finding that order is the hardest part to
estimate~\cite{wada2025hyperbolic} and the over-complete-then-prune approach
of~\cite{lee2025reverse}. Input gain is a search parameter, because drive
depends on level~\cite{sd4562026nameqmatch}.

\paragraph{Initial estimate} Train an encoder on renders of the template chain
with random settings drawn around the factory presets, on DI that may be
shipped (GuitarSet, Guitar-TECHS, Rossi et al.'s improvisations), with the
renders mixed over backing tracks and then separated by the same BS-RoFormer,
so that the encoder sees separation artefacts (the StemFX recipe~\cite{cheng2026stemfx}).
Heads predict the gain class, the effects present and the continuous
parameters, pre-trained on parameter loss~\cite{masuda2021ddsp}, and return
the top $k$ candidates. AFx-Rep (Apache-2.0) is a starting point for the
encoder~\cite{steinmetz2024stito}. Share each dry clip across many chains in
training, or the encoder learns the riff instead of the
tone~\cite{nathanben2026tonerunner}. Before this model exists, the nearest
factory preset by the objective below is the initial estimate.

\paragraph{Refinement} Run CMA-ES from each candidate (population 32--64,
25--50 generations as in ST-ITO) by rendering the pseudo-dry signal through
the offline renderer. The objective has four terms, all after loudness
normalisation:
\begin{itemize}
\item an \emph{aligned spectral} term, MR-STFT and mel distance, for pseudo-dry
signals aligned with the stem (resynthesis, removal), the metric that agrees
best with listeners (Table~\ref{tab:metrics});
\item an \emph{alignment-free} term: long-term spectra per pitch register,
crest factor, compression and high-band flatness statistics, which also work
with the player's DI~\cite{sd4562026nameqmatch};
\item an \emph{effect-embedding} term (AFx-Rep, or a relative encoder such as
RelFx~\cite{liu2026beyond}), with a small weight, because an embedding alone
is an unreliable objective~\cite{long2026tape};
\item a \emph{prior}: a Gaussian mixture over normalised settings of factory
and user presets~\cite{yu2025improving}.
\end{itemize}
From the measured per-sample costs, a 20-second excerpt renders in about
0.4\,s without a capture and 2.3\,s with one on one core, so the ST-ITO budget
of 1,600 renders takes roughly 1.5 to 8 minutes on eight cores (an estimate,
not a measurement). Shorter excerpts in early generations reduce it.

\paragraph{Time-based effects from the notes} The transcription tells where
notes end and where the part rests, which the literature does not exploit. Delay
time: score candidate values tied to the tempo (quarter, dotted eighth,
eighth) and free values by the energy that repeats at that lag after note
offsets. Reverberation: fit the decay of the stem in rests after the last
note and map it to the reverb pedal's settings through a table measured on the
pedal. Modulation: the rate from the modulation spectrum of sustained notes,
or from an LFO extractor~\cite{mitcheltree2023modulation} (MPL-2.0). These
values initialise or fix the corresponding parameters of the search, which
avoids the rugged landscapes of Section~\ref{sec:comparison}.

\paragraph{Linear match} After the search, compute the smoothed difference
between the long-term spectra of the stem and of the rendered chain, and
either fit the equaliser pedal to it or write a minimum-phase impulse response
for the cabinet slot, with a gain limit, as Tone Match
does~\cite{fractalndtonematch}. This absorbs pickup, microphone and mix
equalisation, and it is the step products show to be reliable once gain is
right.

\paragraph{Output} Save each estimate as a user preset with the song and artist
set, so the existing tone suggestion offers it on the Tabs page, keep the
other top candidates, and let the user compare the preset with the stem played
in sync by the tab player. Record the estimate's settings and objective values
in the tab's metadata sidecar for later evaluation.

\paragraph{Captures} When the user has NAM captures on disk, add them as
candidates for the amplifier slot, ranked by rendering the pseudo-dry signal
through each at several input gains~\cite{sd4562026nameqmatch}. TONE3000 can be
queried by metadata (gear type, tone type) for candidates the user then
chooses to download; the catalogue cannot be pre-indexed. Learning a capture
from the stem alone, by unpaired or diffusion-based
fitting~\cite{wright2023adversarial,moliner2025unsupervised}, is a research
direction for tones the chain cannot express, at hours of GPU time per tone.
A neural tone model of any kind could then be converted into a \texttt{.nam}
file by rendering NAM's standard input through it and training a capture on the
pair, so the result plays in the existing NAM pedal (untested).

\subsection{Why this design}

\begin{itemize}
\item \emph{Against direct regression only:} synthetic-to-real gaps and
ambiguous settings; refinement by rendering gained 4.9\,dB in Okita and
Katayose~\cite{okita2026audio}.
\item \emph{Against embedding-only search:} ST-ITO's authors report failure on
guitar tone, and TAPE shows the objective drifting~\cite{steinmetz2024stito,long2026tape};
the aligned pseudo-dry signal gives a spectral term instead.
\item \emph{Against neural tone transfer:} EG-VAE gives the best guitar tone
similarity, but returns audio, has no released code, and cannot run in the
real-time chain~\cite{yeh2026egvae}; it remains useful as a remover or an
encoder if released.
\item \emph{Against linear matching only:} spectral matching cannot set gain or
dynamics, and its own manual requires the gain to be right
first~\cite{fractalndtonematch}.
\item \emph{For the application's own chain:} its parameters are what the
player edits, its renderer is the exact forward model, and training on its
renders gives labels in its own parameter space with licences under our
control.
\end{itemize}

\subsection{Evaluation protocol}
\label{sec:protocol}

\begin{enumerate}
\item \emph{In-chain synthetic.} Held-out DI (Guitar-TECHS P3, GuitarSet
players 03--05) through random template chains, isolated and over backing
tracks after separation. Report parameter error, MR-STFT and MFCC distances
between the true DI rendered through the estimate and the target, and the
embedding distance. The difference between isolated and separated inputs
measures the cost of separation.
\item \emph{Equipment outside the chain.} EGDB-PG renders (CC BY), Guitar-TECHS
real amplifiers, and for research only GOAT and ToneTwist devices: render the
true DI through the estimate and compare with the real output (audio metrics
only).
\item \emph{Real productions.} MoisesDB songs with isolated guitar tracks:
compare estimates from the separated stem with estimates from the isolated
track, and score each against the isolated track with the audio metrics.
\item \emph{Listening.} MUSHRA with webMUSHRA on sets 1 and 2, where the same
performance exists on both sides; tone-similarity opinion scores on set 3.
Report how the objective metrics correlate with the ratings.
\end{enumerate}
Baselines: the current metadata-based suggestion; the nearest factory preset
by the objective; that preset plus the linear match; the full pipeline; and,
as an upper bound, the search with the true DI.

\subsection{Order of work}

\begin{enumerate}
\item The benchmark (sets 1--3) and the baselines, then the template chain with
CMA-ES on resynthesised pseudo-dry signals and the linear match. This needs
no training and measures how far analysis by synthesis goes alone.
\item The excerpt selection and the note-based estimates of delay, reverb and
modulation.
\item The learned initial estimate on separated renders, trained on the DGX
Sparks.
\item Pseudo-dry signals from a removal model trained on the application's
renders, and the player's DI through a pre-chain tap.
\item Capture candidates from the user's library; unpaired capture fitting as
research.
\end{enumerate}
